\section{Sharpness, reductions and further questions}\label{sec:08-boundaries}
\begin{corollary}\label{corollary:08-boundaries-1}
Among standard games with three strategies per player and full actual player
image, six is the minimum player count for nonsplitting, and the only player
count at which nonsplitting can occur. The six-player witness can have binary payoffs
and common interest. The corresponding exact, cardinal and ordinal statements
hold separately.
\end{corollary}
\begin{proof}
For $n=1$ the quotient is trivial; for $n=2$ the standard subgroup already is
a full-image strategy-trivial section. The three-player result was proved
in Theorem\ref{thm:ternary}. The player-count results of the preceding section
exclude all other $n\ne6$. Theorem\ref{thm:six} attains 6. For utility-level
transfers apply the literal normalized-group identities stated after
Theorem\ref{thm:ternary}. The binary witness is unchanged by normalization.
\end{proof}

This is a boundary for \emph{nonsplitting}, not for absence of a full matching:
Example 5.13 already supplies the latter at three players with two strategies.
At three players four strategies are necessary and sufficient for standard
full-image nonsplitting. The lower bound follows from Theorem~\ref{thm:ternary} for three strategies
and Proposition~\ref{prop:binary} for two; with one strategy the full player
image itself gives a section. An explicit binary witness uses strategies $\mathbb Z/4$, players
$\mathbb Z/3$, and $F(d,e)=1$ iff $d<e$ and $d\le1$, using representatives 0 to 3.
Set $u_i(s)=F(s_{i+1}-s_i,s_{i+2}-s_i)$ when $s_i$ is even, and reverse $F$'s
arguments otherwise. Its complete group consists of common translations $T_k$
paired with $p\in S_3$ such that $k\bmod2=\varepsilon(p)$. Every transposition
lift has order 4. The supplementary analytic proof establishes that exact
complete group; the witness is not needed for the six-player lower bound.

A mechanism records the actual represented player projection, rather than
just the abstract group. Structural inequivalence need not imply
irreducibility: the four-strategy witness restricted to common $\{0,1\}$ is
Example 5.13 with $\alpha=\gamma=\delta=0$, $\beta=1$. Its new signed
symmetries need not be inherited from the original game.

By contrast common interest survives exact isomorphisms, rectangular strategy
restrictions, full-payoff clone quotients, and deletion of a player at a fixed
action. No sequence of these operations can turn either common-interest
orientation witness into a non-full Example 5.13, whose majority utilities
$\beta,\gamma$ differ. The six-player code game is clone free: every label
occurs in every coordinate, and changing one coordinate of a winning profile
cannot remain winning because the minimum code distance is 4. These are
precisely scoped reduction statements, not a universal canonical-core theorem.

The most useful next structural question is classification of the represented
six-player obstruction and its payoff-realizable actions, rather than another
increase in enumeration limits. Is every nonsplit ternary action equivalent
under player and local strategy relabelings to the displayed representation?
A group-isomorphism classification alone would not answer this question.
Related problems concern minimal exact realizations, strategy counts beyond
three and reductions admitting utility transformations. A symbolic proof of
the three-player finite lemma would also improve the present account.

The companion fairness project studies utility invariance, assignment models
and ordered versus unordered opponents on independent decision primitives.
No subjective assignment-indifference or outcome-fairness claim is deduced
from the structural theorems here.
